Large language model agents can interpret documents and invoke corporate tools, but they may also treat instructions in external content as legitimate commands. In accounts payable, indirect prompt injection can alter financial fields, redirect beneficiaries, or induce unauthorized critical actions. This study aims to develop and evaluate a protection architecture integrating LLMs with Odoo 19.0 Community in a synthetic vendor-bill workflow. Following Design Science Research and a controlled experiment, it compares a functional reference integration (C0) with defense in depth (C1), combining deterministic policy, least privilege, action-bound approval, idempotency, and auditing. The main experiment specifies 90 benign--adversarial counterfactual pairs, three repetitions, and 1,080 executions with one DeepSeek model; a separate pilot adds 144. Evaluation observes ERP state, security, and utility non-inferiority; latency and cost will be described. Technical proofs cover the document workflow, private C0/C1 execution paths, state capture and restoration, and records preserving failures. These findings do not establish protection effectiveness: operational integration, review, pilot, and protocol freeze remain pending. Scientific conclusions depend on complete data collection. Cross-model replication, ablation, and defense-aware attacks are future work, limiting external validity and attribution to individual layers.
